\documentclass[10pt,a4paper]{amsart}
\usepackage{amsmath,amssymb,amsthm,graphicx}
\usepackage{verbatim}
\newenvironment{reference}{\comment}{\endcomment}
%\newenvironment{reference}{}{}
\newenvironment{slogan}{\comment}{\endcomment}
\newenvironment{history}{\comment}{\endcomment}

% For commutative diagrams we use Xy-pic
\usepackage[all]{xy}

% We use 2cell for 2-commutative diagrams.
\xyoption{2cell}
% Theorem environments.
%
\theoremstyle{plain}
\newtheorem{theorem}[subsection]{Theorem}
\newtheorem{proposition}[subsection]{Proposition}
\newtheorem{lemma}[subsection]{Lemma}

\theoremstyle{definition}
\newtheorem{definition}[subsection]{Definition}
\newtheorem{example}[subsection]{Example}
\newtheorem{exercise}[subsection]{Exercise}
\newtheorem{situation}[subsection]{Situation}

\theoremstyle{remark}
\newtheorem{remark}[subsection]{Remark}
\newtheorem{remarks}[subsection]{Remarks}

\numberwithin{equation}{subsection}

% Macros
%
\def\lim{\mathop{\mathrm{lim}}\nolimits}
\def\colim{\mathop{\mathrm{colim}}\nolimits}
\def\Spec{\mathop{\mathrm{Spec}}}
\def\Hom{\mathop{\mathrm{Hom}}\nolimits}
\def\Ext{\mathop{\mathrm{Ext}}\nolimits}
\def\SheafHom{\mathop{\mathcal{H}\!\mathit{om}}\nolimits}
\def\SheafExt{\mathop{\mathcal{E}\!\mathit{xt}}\nolimits}
\def\Sch{\mathit{Sch}}
\def\Mor{\mathop{\mathrm{Mor}}\nolimits}
\def\Ob{\mathop{\mathrm{Ob}}\nolimits}
\def\Sh{\mathop{\mathit{Sh}}\nolimits}
\def\NL{\mathop{N\!L}\nolimits}
\def\CH{\mathop{\mathrm{CH}}\nolimits}
\def\proetale{{pro\text{-}\acute{e}tale}}
\def\etale{{\acute{e}tale}}
\def\QCoh{\mathit{QCoh}}
\def\Ker{\mathop{\mathrm{Ker}}}
\def\Im{\mathop{\mathrm{Im}}}
\def\Coker{\mathop{\mathrm{Coker}}}
\def\Coim{\mathop{\mathrm{Coim}}}

% Boxtimes
%
\DeclareMathSymbol{\boxtimes}{\mathbin}{AMSa}{"02}

%
% Macros for moduli stacks/spaces
%
\def\QCohstack{\mathcal{QC}\!\mathit{oh}}
\def\Cohstack{\mathcal{C}\!\mathit{oh}}
\def\Spacesstack{\mathcal{S}\!\mathit{paces}}
\def\Quotfunctor{\mathrm{Quot}}
\def\Hilbfunctor{\mathrm{Hilb}}
\def\Curvesstack{\mathcal{C}\!\mathit{urves}}
\def\Polarizedstack{\mathcal{P}\!\mathit{olarized}}
\def\Complexesstack{\mathcal{C}\!\mathit{omplexes}}
% \Pic is the operator that assigns to X its picard group, usage \Pic(X)
% \Picardstack_{X/B} denotes the Picard stack of X over B
% \Picardfunctor_{X/B} denotes the Picard functor of X over B
\def\Pic{\mathop{\mathrm{Pic}}\nolimits}
\def\Picardstack{\mathcal{P}\!\mathit{ic}}
\def\Picardfunctor{\mathrm{Pic}}
\def\Deformationcategory{\mathcal{D}\!\mathit{ef}}
\usepackage{cleveref}
\crefname{originalreference}{original reference}{original references}
\Crefname{originalreference}{Original reference}{Original references}
\makeatletter
\newlabel{item-construction-associated-cofibered-groupoid}{{06GN}{source}{}{sourcepointer-1}{}}
\newlabel{item-construction-associated-cofibered-groupoid@cref}{{[originalreference][0][]06GN}{[1][0][]source}{}{}{}}
\makeatother
\crefname{originalreference}{original reference}{original references}
\Crefname{originalreference}{Original reference}{Original references}
\makeatletter
\newlabel{remarks-cofibered-groupoids}{{06GK}{source}{}{sourcepointer-2}{}}
\newlabel{remarks-cofibered-groupoids@cref}{{[originalreference][0][]06GK}{[1][0][]source}{}{}{}}
\makeatother
\title{Smoothness of a groupoid in functors is equivalent to smoothness of its source map and of its quotient morphism}
\author{The Stacks Project Authors}
\date{stacks-project:formal-defos.tex:local-snapshot-20261012; GNU FDL 1.2 or later}
\begin{document}
\maketitle
\noindent\textit{Editorial note.} All delivered mathematical text is the author's own, copied from the pinned Stacks Project chapter file formal-defos.tex; no statement or proof step was rewritten, reordered or removed.\par\medskip
\noindent\textit{Editorial note.} The counted proof is the author's complete argument: the comparison of the source and target maps through the inverse morphism of the groupoid, the explicit unwinding of the two smoothness conditions, and the translation between them using the inverse morphism.\par\medskip
\noindent\textit{Editorial note.} The label lemma-smooth-quotient-morphism names the counted statement, and the proof environment immediately following it in the source is delivered verbatim as its proof.\par\medskip
\noindent\textit{Editorial note.} The counted statement is carried by the author's own lemma environment in the chapter source.\par\medskip
\noindent\textit{Editorial note.} Pre-open gate: python operations/agent-authoring/tools/precard\_gate.py --source formal-defos.tex --label lemma-smooth-quotient-morphism --start 6656 --end 6699 --local-id 28022 returned RESULT=OK (L28022 free, occupancy FREE, public interval clear). Candidate taken from sources/stacks-formal-defos/candidates-halfA.json.\par\medskip
\noindent\textit{Editorial note.} The statement and proof are copied from the pinned chapter file; the cited inputs are kept as external pointers with their canonical Stacks tags.\par\medskip
\noindent\textit{Allowed standard prerequisites (editorial).}\begin{itemize}
\item Inverse morphism of the groupoid and the identity section. Same chapter, the defining data of a groupoid in functors, in particular the inverse morphism i : R -> R and the identity e : U -> R. Supplies the identities t = s composed with i and the elements used to translate between the two lifting conditions.
\item Description of morphisms of the quotient groupoid. Same chapter, the remark describing the objects and morphisms of the quotient [U/R] as pairs with the compatibility conditions. Supplies the translation of a morphism (f, a) of the quotient into the data of ring maps and elements used in the unwinding.
\end{itemize}\medskip
\ifdefined\hypertarget\hypertarget{sourcepointer-1}{}\fi
\noindent\textit{Source reference.} Original reference 06GN: sources/stacks-formal-defos/pinned-original-source.tex, source lines 1432--1432. Prerequisite pointer invoked at proof line(s) 6690; it names a standard fact of the same chapter that the counted argument uses, and every step of the counted proof built on it is retained verbatim in the delivered body.\par\medskip
\ifdefined\hypertarget\hypertarget{sourcepointer-2}{}\fi
\noindent\textit{Source reference.} Original reference 06GK: sources/stacks-formal-defos/pinned-original-source.tex, source lines 1381--1381. Prerequisite pointer invoked at proof line(s) 6689; it names a standard fact of the same chapter that the counted argument uses, and every step of the counted proof built on it is retained verbatim in the delivered body.\par\medskip
\begin{lemma}
\label{lemma-smooth-quotient-morphism}
Let $(U, R, s, t, c)$ be a groupoid in functors on $\mathcal{C}_\Lambda$.
The following are equivalent:
\begin{enumerate}
\item The groupoid in functors $(U, R, s, t, c)$ is smooth.
\item The morphism $s : R \to U$ is smooth.
\item The morphism $t : R \to U$ is smooth.
\item The quotient morphism $U \to [U/R]$ is smooth.
\end{enumerate}
\end{lemma}

\begin{proof}
Statement (2) is equivalent to (3) since the inverse $i: R \to R$ of
$(U, R, s, t, c)$ is an isomorphism and $t = s \circ i$. By definition (1) is
equivalent to (2) and (3) together, hence it is equivalent to either of them
individually.

\medskip\noindent
Finally we prove (2) is equivalent to (4). Unwinding the definitions:
\begin{enumerate}
\item[(2)] Smoothness of $s: R \to U$ amounts to the following
condition: If $f: B \to A$ is a surjective ring map in
$\mathcal{C}_\Lambda$, $a \in R(A)$, and $y \in U(B)$ such that
$s(a) = U(f)(y)$, then there exists $a' \in R(B)$ such that
$R(f)(a') = a$ and $s(a') = y$.
\item[(4)] Smoothness of $U \to [U/R]$ amounts to the following
condition: If $f: B \to A$ is a surjective ring map in
$\mathcal{C}_\Lambda$ and $(f, a) : (B, y) \to (A, x)$
is a morphism of $[U/R]$, then there exists $x' \in U(B)$ and
$b \in R(B)$ with $s(b) = x'$, $t(b) = y$
such that $c(a, R(f)(b)) = e(x)$.  Here $e : U \to R$ denotes the
identity and the notation $(f, a)$ is as in Remarks
\ref{remarks-cofibered-groupoids}
(\ref{item-construction-associated-cofibered-groupoid});
in particular $a \in R(A)$ with $s(a) = U(f)(y)$ and $t(a) = x$.
\end{enumerate}
If (4) holds and $f, a, y$ as in (2) are given, let $x = t(a)$ so that
we have a morphism $(f, a): (B, y) \to (A, x)$.
Then (4) produces $x'$ and $b$, and $a' = i(b)$ satisfies the requirements
of (2). Conversely, assume (2) holds and let $(f, a): (B, y) \to (A, x)$
as in (4) be given. Then (2) produces $a' \in R(B)$, and $x' = t(a')$ and
$b = i(a')$ satisfy the requirements of (4).
\end{proof}

\end{document}
